\section{国民年金基金}

\subsection{制度の概要}

\begin{fillblank}[年金1 2019 問題1(4)\textcircled{1}]
  2019 年4 月1 日に（　　A　　）の（　　B　　）型国民年金基金と全国47 都道府県の（　　C　　）型国
民年金基金が統合され、全国国民年金基金が設立されている。
\end{fillblank}

\begin{answerbox}
  \begin{enumerate}[A:]
    \item 22
    \item 職能
    \item 地域
  \end{enumerate}
\end{answerbox}

\begin{fillblank}[年金1 H29 問題1(3)]
  国民年金基金および国民年金基金連合会について、次の空欄に当てはまる適
切な語句を記入しなさい。

\begin{itemize}
  \item 国民年金基金は国民年金の（　　A　　）を代行している。このため、国民年金基金に加入している
者は、国民年金基金の保険料とは別に国民年金の（　　A　　）の保険料を納付することができない。

  \item 国民年金基金の給付の型は、全員が加入する１口目（給付の型は終身年金が２種類）と希望
に応じて加入する２口目以降（給付の型は終身年金が２種類、確定年金が（　　B　　）種類）があり、
加入後も加入口数を増減することができる。ただし、確定年金の年金額は、終身年金の年金額
（１口目を含めた額）を超えることはできない。

  \item 国民年金基金の掛金額は、 （　　C　　） 、性別、加入する給付の型および口数によって決まる。
  
  \item 国民年金基金の加入員資格を 60 歳到達前に喪失した者で加入員期間が（　　D　　）年未満の者は、
国民年金基金連合会に年金資産を移し、加入していた国民年金基金に代わって国民年金基金連
合会から給付が支給される。

  \item 確定拠出年金法第 59 条の規定により、国民年金基金連合会は、少なくとも 5 年ごとに、個人
型年金加入者数の動向、企業型年金の実施の状況、国民生活の動向等を勘案し、（　　E　　）の内容
について再検討を加え、必要があると認めるときは（　　E　　）を変更しなければならない。
\end{itemize}

\end{fillblank}

\begin{answerbox}
  \begin{enumerate}[A:]
    \item 付加年金
    \item 5
    \item 加入時の年齢
    \item 15
    \item 個人型年金規約
  \end{enumerate}
\end{answerbox}

\begin{fillblank}[数理人試験 2019 問題1 設問6 A,E,F]
  次は、国民年金基金に関する記述である。
  \begin{itemize}
    \item 2019年4月に全国47都道府県の地域型国民年金基金と22の職能型国民年金基金が合併し、
（　　A　　）となった。
    \item 国民年金の第1号被保険者であっても、（　　B　　）は国民年金基金に加入できない。
    \item 国民年金基金の掛金は予定利率（　　C　　）で計算されている。
  \end{itemize}

\end{fillblank}

\begin{answerbox}
  \begin{enumerate}[A:]
    \item 全国国民年金基金
    \item 国民年金の保険料を免除されている者（または、農業者年金の被保険者）
    \item 1.5％
  \end{enumerate}
\end{answerbox}


\begin{fillblank}[数理人試験 2024 問題1 設問5 A,B,E]
  次は、国民年金基金に関する記述である。
  \begin{itemize}
    \item 国民年金基金は、「全国国民年金基金」と （　　A　　）つの職種別に設立された「職能型国民年金
基金」の2種類がある。
\item 職能型国民年金基金に加入できるのは、基金ごとに定められた事業または業務に従事する国民年金の
第（　　B　　）号被保険者である。
\item 国民年金基金の加入員資格を60歳到達前に喪失した者で加入員期間が（　　C　　）年未満の者につい
ては、国民年金基金連合会に年金資産を移し、加入していた国民年金基金に代わって国民年金基金連合
会から支給される。
  \end{itemize}
\end{fillblank}

\begin{answerbox}
  \begin{enumerate}[A: ]
    \item 3
    \item 1
    \item 15
  \end{enumerate}
\end{answerbox}


\newpage

\subsection{掛金}

\begin{fillblank}[年金1 2025 問題1(6)ア]
  国民年金基金の現行の掛金額の上限は月額 （　　A　　） 万円（個人型確定拠出年金の掛金との
合算）であるが、令和6 年（2024 年）12 月27 日に閣議決定された「令和７年度税制改正の大
綱」によると、月額 （　　B　　） 万円に引き上げることとされた。
\end{fillblank}

\begin{answerbox}
    \begin{enumerate}[A: ]
    \item 6.8
    \item 7.5
  \end{enumerate}
\end{answerbox}

\begin{fillblank}[自作 2026/10/31]
  国民年金の保険料を免除（一部免除・学生納付特例・納付猶予を含む。）されていた者が、
  直近10年以内の免除された全期間分の保険料を追納した場合は、
  国民年金の保険料が免除されていたことにより国民年金基金に加入できなかった期間に相当する期間
  （（　　A　　）年を限度とする。）に限り、掛金を月額（　　B　　）円まで
  （令和8年12月分掛金からは（　　C　　）円まで）納めることができる。
\end{fillblank}

\begin{answerbox}
    \begin{enumerate}[A: ]
    \item 5
    \item 102,000
    \item 112,500
  \end{enumerate} 
\end{answerbox}

\begin{explanationbox}
  令和7年税制改正で追納する時の上限も引き上がることになった。

11.25万円を忘れた場合の逆算法：$10.2万円\times7.5万円÷6.8万円=11.25万円$
\end{explanationbox}

\begin{fillblank}[数理人試験 2019 問題1 設問6 B,C,D 改題]
  次は、国民年金基金に関する記述である。
  \begin{itemize}
    \item 国民年金基金の掛金の上限は、令和８年12月分掛金を境に原則月額（　　A　　）円から
    月額（　　B　　）円に引き上がる。ただし、（　　C　　）にも加入
している場合は、その掛金と合わせて（　　B　　）円以内となる。
    \item 国民年金基金の掛金は（　　D　　）の対象となる。
  \end{itemize}
\end{fillblank}

\begin{answerbox}
  \begin{enumerate}[A: ]
    \item 68,000
    \item 75,000
    \item 個人型確定拠出年金
    \item 社会保険料控除
  \end{enumerate}
\end{answerbox}

\begin{shortanswer}[年金1 2022 問題2(4)イ]
  国民年金基金の掛金について、月額の算定方法、上限および払込期間を簡潔に入力しなさい。
\end{shortanswer}

\begin{answerbox}
  ＜月額の算定方法＞

掛金の月額は、加入時年齢・性別・選択した給付の型と加入口数に応じて決まり、加入時年齢・
性別・給付の型ごとに1 口あたりの掛金月額を定めた掛金月額表により掛金が算定される。

＜上限＞

掛金の上限は、令和8年12月分掛金から月額75,000 円であり、個人型年金にも加入している場合は、
その掛金と合わせて月額75,000 円となる。

＜払込期間＞

掛金の払込期間は、加入時から60 歳に達する日の属する月（60 歳以上で加入した場合は65 歳
に達する日の属する月または国民年金の任意加入被保険者資格の喪失予定年月）の前月まで。
\end{answerbox}

\begin{shortanswer}[年金1 2024 問題2(4)]
  国民年金基金では掛金の上限額は 68,000 円（令和8年12月分掛金からは75,000円）とされているが、
  国民年金の保険料を免除されて
いた者について掛金の額の上限が増加する特例がある。当該特例の適用条件、適用期間およ
び掛金の上限額について簡潔に入力しなさい。 （250 字以内）
\end{shortanswer}

\begin{answerbox}
  ＜適用条件＞

加入員が保険料免除期間における保険料の全部の追納を行った場合又は当該保険料の全部につ
き追納を行った国民年金の被保険者が加入員となった場合

＜適用期間および上限額＞

追納した期間（５年間が上限）につき、各月102,000 円が上限額となる（令和8年12月1日からは112,500円が上限額となる）
\end{answerbox}

\newpage

\subsection{給付}

\begin{fillblank}[数理人試験 2025 問題1 設問5]
  次は、国民年金基金に関する記述である。

  \begin{itemize}
    \item 国民年金基金の掛金について
    \begin{itemize}
      \item 4月から翌年3月までの1年分の掛金を前納すると掛金が割引され、掛金前納額は次の算式の通りと
なる。
  \begin{flalign*}
        掛金前納額 = １か月の基金掛金額 \times（　　\text{A}　　）
  \end{flalign*}
      \item 掛金月額の上限は68,000円（令和8年12月分掛金からは75,000円）であるが、国民年金の保険料を免除（一部免除・学生納付特例・納付猶
予を含む。）されていた者が、直近10年以内の免除された全期間分の保険料を追納した場合は、国
民年金の保険料が免除されていたことにより国民年金基金に加入できなかった期間に相当する期間
（（　　B　　）年を限度とする。）に限り、掛金を月額102,000円（令和8年12月分掛金からは月額112,500円）まで納めることができる。
    \end{itemize}

    \item 国民年金基金の遺族一時金について
    \begin{itemize}
      \item 保証期間のある終身年金（　　C　　）と確定年金\ajRoman{1}型、\ajRoman{2}型、\ajRoman{3}型、\ajRoman{4}型、\ajRoman{5}型に加入している者
が、年金を受け取る前、又は保証期間中に死亡した場合、遺族に一時金が支払われる。
      \item 保証期間のない終身年金（　　D　　）のみに加入している場合でも年金を受給する前に死亡した場
合、（　　E　　）万円の一時金が遺族に支払われる。
      \item 税制に関して、遺族一時金は（　　F　　）である。
    \end{itemize}
  \end{itemize}

\end{fillblank}

\begin{answerbox}
  \begin{enumerate}[A: ]
    \item 11.9
    \item 5
    \item A型
    \item B型
    \item 1
    \item 非課税
  \end{enumerate}
\end{answerbox}

\begin{explanationbox}
  A: 4月から翌年3月までの1年分の掛金を前納すると0.1ヶ月分の掛金が割引される。
\end{explanationbox}

\begin{fillblank}[数理人試験 2024 問題1 設問5 C,D,F]
  次は、国民年金基金に関する記述である。
  \begin{itemize}
    \item 国民年金基金の付加型である2口目以降の老齢年金の給付の型として、保証期間が10年のものは
（ 　　A　　）種類、保証期間が15年のものは（　　B　　）種類ある。
\item 国民年金基金の加入員資格を60歳到達前に喪失した者で加入員期間が（　　C　　）年未満の者につい
ては、国民年金基金連合会に年金資産を移し、加入していた国民年金基金に代わって国民年金基金連合
会から支給される。
  \end{itemize}
\end{fillblank}

\begin{answerbox}
\begin{enumerate}[A: ]
  \item 2
  \item 3
  \item 15
\end{enumerate}
\end{answerbox}

\begin{shortanswer}[年金1 2022 問題2(4)ア]
  国民年金基金の老齢年金について、給付の型およびその選択方法を簡潔に入力しなさい。
\end{shortanswer}

\begin{answerbox}
  ＜給付の型＞

終身年金のA 型・B 型と、確定年金の\ajRoman{1}型・\ajRoman{2}型・\ajRoman{3}型・\ajRoman{4}型・\ajRoman{5}型の7 種類がある。

＜選択方法＞

1 口目は終身年金 A 型・B 型のいずれかを選択し、2 口目以降は、終身年金 A 型・B 型、確定
年金\ajRoman{1}型・\ajRoman{2}型・\ajRoman{3}型・\ajRoman{4}型・\ajRoman{5}の7 種類の中から自由に組み合わせて選択できる。
ただし、2 口目以降の選択に当たっては、確定年金の年金額が終身年金の年金額（1 口目を含め
た額）を超えるような選択はできない。
\end{answerbox}

\newpage